\chapter{谱混合、质量轴线形与同位素贡献}
\section{碎片响应矩阵}
电子轰击等电离过程可以使一个中性分子形成多个离子碎片。有限监测通道中的单组分指纹记为 $s_j\in\R^M$。在线性响应区间，混合物满足
\begin{equation}
 y_m(t)=\sum_{j=1}^J S_{mj}z_j(t)+b_m(t)+\varepsilon_m(t),
 \qquad Y=SZ+B_0+E.\label{eq:mix}
\end{equation}
$S$ 的第 $j$ 列是组分 $j$ 在本机条件下的相对通道响应。若选择参考通道 $m_j$，可令 $S_{m_jj}=1$，把绝对响应系数留在 $z_j$ 或后述 $\kappa_j$ 中。不同列的归一化方式必须一致地记录，不能在拟合后无痕改变。

NIST Chemistry WebBook 的谱可作为候选结构与碎片模式的参考\cite{nistweb}。相对指纹不等于当前离子源能量、质量透过率、电子倍增器增益和采样条件下的绝对灵敏度。数据库谱可以帮助定义假设，却不能替代本机标准品与混合物验证。
\fig{02_response}{.79}{本报告构造的七通道、三组分响应矩阵。A、B、C 是数学组分标签，不表示已鉴定的化合物；系数不是 NIST 谱库数据。}

\section{质量轴上的卷积}
设理想质量分布为 $a(\mu)$，仪器线形为 $g(m,\mu)$，则观测为
\begin{equation}
 p(m)=\int g(m,\mu)a(\mu)\,\mathrm d\mu+\varepsilon(m).
\end{equation}
只有在所考察窄窗口内线形近似只依赖差值时，才能写成
\begin{equation}
 g(m,\mu)=g(m-\mu),\qquad p=g*a+\varepsilon.\label{eq:massconv}
\end{equation}
宽质量范围内的分辨率变化、峰形不对称变化、质量校准漂移都可能使此近似失效。软件的质量轴路径要求均匀 \code{mz_grid}、奇数长度的居中非负线形以及单位离散和；核中心与质量格点一一对应，避免半格偏移。

式~\eqref{eq:massconv} 的恢复对象仍是一个分辨率受限的质量分布。它不自动完成多电荷包络反演、同位素簇分组、串联质谱库检索或分子结构鉴定。那些任务需要额外离子化与化学模型，未在本版 MCP 接口中实现。

\fig{07_mass_profile}{.93}{窄质量窗口中的一维线形恢复。观测与真值、恢复曲线均来自合成数据；所用线形在此窗口内固定。图像去模糊不属于本例。}

\section{以二氧化碳为例推导同位素比}
设 $p_{12},p_{13}$ 是碳同位素丰度，$p_{16},p_{17},p_{18}$ 是氧同位素丰度。只考虑单电荷分子离子，且暂不考虑质量歧视，则以 $^{12}\mathrm C{}^{16}\mathrm O_2$ 为基准，质量 45 的主要分子组合为
\begin{equation}
^{13}\mathrm C{}^{16}\mathrm O_2,
\qquad ^{12}\mathrm C{}^{16}\mathrm O{}^{17}\mathrm O.
\end{equation}
两个氧位点给出组合因子 2。因此
\begin{equation}
\frac{I_{45}}{I_{44}}
=\frac{p_{13}}{p_{12}}+2\frac{p_{17}}{p_{16}}.
\end{equation}
质量 46 的主要组合为一个 $^{18}\mathrm O$、两个 $^{17}\mathrm O$，以及 $^{13}\mathrm C$ 与一个 $^{17}\mathrm O$，故
\begin{equation}
\frac{I_{46}}{I_{44}}
=2\frac{p_{18}}{p_{16}}
+\left(\frac{p_{17}}{p_{16}}\right)^2
+2\frac{p_{13}}{p_{12}}\frac{p_{17}}{p_{16}}.\label{eq:isotope46}
\end{equation}
采用 NIST 表中代表性组成，可得到约 $1.158\%$ 和 $0.412\%$\cite{nistC,nistO}。天然材料同位素组成存在变化范围\cite{ciaaw}；仪器质量歧视、碎片、其它分子的重叠和背景也会改变实测比例。因此这些数值只是物理检查先验，不是可以无条件套用的背景扣除系数。

\section{同位素诊断与产物鉴定的区别}
在 CO$_2$ 占优的背景中，45 和 46 通道与 44 共同变化是合理的先验现象。若先在可信空白段建立参考通道、压力和缓慢漂移模型，再在反应段检查残差，可以判断是否存在超出该背景模型的附加信号。但附加信号仍可能来自未建模的系统变化。

因此，“45 有残差”与“产生了乙醇”不是等价命题。判断候选产物还需要独立碎片、标准品、空白、标记实验或其它分析方法相互支持。本版的 \code{ms_background} 输出带符号残差、回归条件和子数据集，不输出物种确认标签。

\section{谱可辨识性的初步检查}
若 $\operatorname{rank}(S)<J$，仅利用瞬时谱无法唯一分离 $J$ 个组分；即使满秩，几乎共线的列也会把小噪声放大。为消除列幅值对诊断的干扰，先令
\begin{equation}
\widetilde s_j=s_j/\|s_j\|_2,\qquad
\mu=\max_{i\ne j}|\widetilde s_i^T\widetilde s_j|.
\end{equation}
软件报告归一化奇异值、条件数、最大相干性和近共线列对。该诊断只考察 $S$；它不是完整谱--时间联合算子的可辨识性定理。不同且已知的时间核有时能补充信息，未知核与背景也可能进一步削弱信息。
